\subsection{LIE ALGEBRA OF A LIE GROUP GERM}

In this no. \((G, e, \theta, m)\) denotes a Lie group germ. A large part of the results of the § are still true with the same proof. We shall review those which we shall find useful.

18.1. Let \(\Omega\) be the set of definition of \(m\) . Let \((g, g') \in \Omega\) , \(t \in \mathrm{T}_g^{(\alpha)}(\mathrm{G})\) , \(t' \in \mathrm{T}_{g'}^{(\alpha)}(\mathrm{G})\) . As in no. 1, the convolution product of \(t\) and \(t'\) , denoted by \(t * t'\) , is the image of \(t \otimes t'\) under \(m\) . We write \(\mathrm{U}(\mathrm{G}) = \mathrm{T}_e^{(\alpha)}(\mathrm{G})\) , \(\mathrm{U}_s(\mathrm{G}) = \mathrm{T}_e^{(s)}(\mathrm{G})\) , \(\mathrm{U}^+(\mathrm{G}) = \mathrm{T}_e^{(\alpha)+}(\mathrm{G})\) , \(\mathrm{U}_s^+(\mathrm{G}) = \mathrm{T}_e^{(s)+}(\mathrm{G})\) . For \(t\) , \(t'\) in \(\mathrm{U}(\mathrm{G})\) , \(t * t'\) is defined and belongs to \(\mathrm{U}(\mathrm{G})\) . With the convolution product, \(\mathrm{U}(\mathrm{G})\) is an associative algebra with unit element \(\varepsilon_e\) , filtered by the \(\mathrm{U}_s(\mathrm{G})\) . The canonical isomorphism \(i_{\mathrm{G}, e}\) of gr \(\mathrm{U}(\mathrm{G})\) onto \(\mathrm{TS}(\mathrm{T}_e(\mathrm{G}))\) is an algebra isomorphism.

18.2. Let G, H be Lie group germs and \(\phi: G \to H\) a morphism. If \(t \in \mathrm{U}(\mathrm{G})\) , the image \(\mathrm{U}(\phi)(t)\) of t under \(\phi_{*}\) is an element of U(H) and \(\mathrm{U}(\phi)\) is a morphism of the algebra U(G) into the algebra U(H). The mapping \(\theta: x \mapsto x^{-1}\) of G into G defines a mapping \(t \mapsto t^{\vee}\) of U(G) into U(G). For \(t, t'\) in U(G), the product \(t * t'\) evaluated relative to \(G^{\vee}\) is equal to the product \(t' * t\) evaluated relative to G and \((t * t')^{\vee} = t^{\prime\vee} * t^\vee\) . Then \(\mathrm{U}(\phi)(t^{\vee}) = (\mathrm{U}(\phi)t)^{\vee}\) . If \(G_{1}, \ldots, G_{n}\) are Lie group germs and \(G = G_{1} \times \cdots \times G_{n}\) , the canonical isomorphism of \(\mathrm{U}(G_{1}) \otimes \cdots \otimes \mathrm{U}(G_{n})\) onto U(G) is an algebra isomorphism;

for \(t_{1},\ldots,t_{n}\) in U(G), \((t_{1}\otimes\cdots\otimes t_{n})^{\vee}=t_{1}^{\vee}\otimes\cdots\otimes t_{n}^{\vee}\) . Let H be a Lie subgroup germ of G and i: H→G the canonical injection. Then U(i) is an injective homomorphism of the algebra U(H) into U(G) and

\[
\mathrm{U} (i) \left(t ^ {\vee}\right) = \left(\mathrm{U} (i) (t)\right) ^ {\vee}
\]

for all \(t \in \mathrm{U}(\mathrm{H})\) . With the convolution product and the coproduct defined by the manifold structure on G, \(\mathrm{U}(\mathrm{G})\) is a bigebra and \(\mathrm{U}(\phi)\) is a bigebra morphism.

18.3. Let G be a Lie group germ, X a manifold of class \(C^{r}\) and \(\psi\) a law chunk of left operation of class \(C^{r}\) of G on X. Let \(\Omega\) be the set of definition of \(\psi\) . If \(t \in T_{g}^{(s)}(G)\) , \(u \in T_{x}^{(s^{\prime\prime})}(X)\) , \((g, x) \in \Omega\) and \(s + s^{\prime} \leqslant r\) , let \(t * u\) denote the image of \(t \otimes u\) under \(\psi_{*}\) . Let \(t \in T_{g}^{(s)}(G)\) , \(t^{\prime} \in T_{g'}^{(s^{\prime})}(G)\) , \(u \in T_{x}^{(s^{\prime})}(X)\) ; if \(s + s^{\prime} + s^{\prime\prime} \leqslant r\) and \(gg^{\prime}\) , \((gg^{\prime})x\) , \(g^{\prime}x\) , \(g(g^{\prime}x)\) are defined, then

\[
(t * t ^ {\prime}) * u = t * (t ^ {\prime} * u).
\]

Let \(x_0 \in \mathbf{X}\) and \(\rho(x_0)\) be the mapping \(g \mapsto gx_0\) , which is defined in an open neighbourhood of \(e\) . If \(t \in \mathrm{U}_r(\mathrm{G})\) , then \(\rho(x_0)_*t = t * \varepsilon_{x_0}\) . Here and in the rest of this no., we shall leave to the reader the task of translating the results for law chunks of right operation.

18.4. Preserving the notation of 18.3, let \(t \in \mathrm{U}_s(\mathbf{G})\) with \(s \leqslant r\) . Let \(f\) be a function of class \(\mathbf{C}^r\) on \(\mathbf{X}\) with values in a Hausdorff polynormed space. Let \(t * f\) denote the function on \(\mathbf{X}\) defined by

\[
\begin{array}{r l} (t * f) (x) & = \langle t, g \mapsto f (\psi (\theta (g), x)) \rangle \\ & = \langle t ^ {\vee}, f \circ \rho (x) \rangle = \langle \rho (x) _ {*} (t ^ {\vee}), f \rangle = \langle t ^ {\vee} * \varepsilon_ {x}, f \rangle . \end{array}
\]

If \(t \in \mathrm{U}_s(\mathrm{G})\) , \(t' \in \mathrm{U}_{s'}(\mathrm{G})\) and \(s + s' \leqslant r\) , then \(\langle t', t * f \rangle = \langle t^\vee * t', f \rangle\) and \((t * t') * f = t * (t' * f)\) . Let \(t \in \mathrm{U}_s(\mathrm{G})\) , \(f\) and \(f'\) be functions of class \(C^r\) on X with values in Hausdorff polynormed spaces \(F, F'\) and \((u, u') \mapsto u. u'\) be a continuous bilinear mapping of \(F \times F'\) into a Hausdorff polynormed space; let \(\underline{n}\) .

\(\sum_{i=1}^{n} t_i \otimes t_i'\) be the image of \(t\) under the coproduct; if \(s \leqslant r\) , then

\[
t * (f f ^ {\prime}) = \sum_ {i = 1} ^ {n} (t _ {i} * f) (t _ {i} ^ {\prime} * f ^ {\prime}).
\]

18.5. Preserving the notation of 18.3, let \(t \in \mathrm{U}_s(\mathrm{G})\) with \(s \leqslant r\) . The mapping \(x \to t * \varepsilon_x\) is called the field of point distributions defined by \(t\) and the law chunk of operation and is sometimes denoted by \(\mathrm{D}_t^{\psi}\) or \(\mathrm{D}_t\) . If \(f: \mathrm{X} \to \mathrm{F}\) is a function of class \(C^r\) , the function \(t^{\vee} * f\) on \(\mathrm{X}\) is also denoted by \(\mathrm{D}_t f\); it is of class \(C^{r - s}\) if \(s < \infty\) . If \(t \in \mathrm{U}_s(\mathrm{G})\) , \(t' \in \mathrm{U}_{s'}(\mathrm{G})\) and \(s + s' \leqslant r\) , then \(\mathrm{D}_{t,t'} f = \mathrm{D}_{t'}(\mathrm{D}_t f)\) . If \(G\) and \(X\) are finite-dimensional, \(D_t\) is a differential operator on \(X\) of order \(\leqslant s\) and of class \(C^{r - s}\) (if \(s < \infty\) ). The function \(D_t f\) is then the transform of \(f\) under this differential operator.

18.6. Let G be a Lie group germ and \(t \in \mathrm{U}(\mathrm{G})\) . \(L_{t}\) denotes the field of point distributions \(g \mapsto \varepsilon_{g} * t\) on G and \(R_{t}\) the field of point distributions \(g \mapsto t * \varepsilon_{g}\) on G. If \(f \in \mathcal{C}^{\omega}(\mathrm{G}, \mathrm{F})\) , then \(\mathrm{L}_{t} f \in \mathcal{C}^{\omega}(\mathrm{G}, \mathrm{F})\) and \(\mathrm{R}_{t} f \in \mathcal{C}^{\omega}(\mathrm{G}, \mathrm{F})\) . For \(t, t'\) in \(\mathrm{U}(\mathrm{G})\) , \(L_{t* t'} = L_t \circ L_{t'}, R_{t* t'} = R_{t'} \circ R_t, L_t \circ R_{t'} = R_{t'} \circ L_t, \theta(L_t) = R_t^\vee\) .

18.7. As \(T_{e}(G)\) is the set of primitive elements of U(G),

\[
[ \mathrm{T} _ {e} (\mathrm{G}), \mathrm{T} _ {e} (\mathrm{G}) ] \subset \mathrm{T} _ {e} (\mathrm{G}).
\]

The normable space \(\mathrm{T}_{e}(\mathrm{G})\) , together with the bracket, is a normable Lie algebra, called the normable Lie algebra of G (or Lie algebra of G) and denoted by \(\mathrm{L}(\mathrm{G})\) . Let \(\mathrm{E}(\mathrm{G})\) be the enveloping algebra of \(\mathrm{L}(\mathrm{G})\) . The canonical injection of \(\mathrm{L}(\mathrm{G})\) into \(\mathrm{U}(\mathrm{G})\) defines a homomorphism \(\eta\) of the algebra \(\mathrm{E}(\mathrm{G})\) into the algebra \(\mathrm{U}(\mathrm{G})\) ; if K is of characteristic 0, \(\eta\) is a bigebra isomorphism, by means of which \(\mathrm{U}(\mathrm{G})\) is identified with \(\mathrm{E}(\mathrm{G})\) . Using the notation of 18.3, for all \(a \in \mathrm{L}(\mathrm{G})\) , let \(D_{a}\) be the field of point distributions defined by a on X. The mapping \((a, x) \mapsto D_{a}(x)\) is a morphism of class \(C^{r-1}\) of the trivial vector bundle \(\mathrm{L}(\mathrm{G}) \times \mathrm{X}\) into the vector bundle \(\mathrm{T}(\mathrm{X})\) . Let I be an open subset of K containing 0 and \(\gamma: I \to G\) a mapping of class \(C^{r}\) such that \(\gamma(0) = e\) . Let \(a = T_{0}(\gamma) 1 \in L(\mathrm{G})\) . If f: X → F is a function of class \(C^{r}\) , then

\[
\left(\mathrm{D} _ {a} f\right) (x) = \lim _ {k \in \mathrm{K} ^ {*}, k \rightarrow 0} k ^ {- 1} \left(f (\gamma (k) x) - f (x)\right).
\]

If \(r \geqslant 2\) , the mapping \(a \mapsto \mathrm{D}_a\) is a law of left infinitesimal operation of class \(\mathrm{C}^{r-1}\) of \(\mathrm{L}(\mathrm{G})\) on \(\mathrm{X}\) .

18.8. Let G and H be Lie group germs and \(\phi\) a morphism of G into H. The restriction of \(\mathrm{U}(\phi)\) to \(\mathrm{L}(\mathrm{G})\) , which is just \(\mathrm{T}_{e}(\phi)\) , is a continuous morphism of \(\mathrm{L}(\mathrm{G})\) into \(\mathrm{L}(\mathrm{H})\) , which we denote by \(\mathrm{L}(\phi)\) . If \(\psi\) is a morphism of H into a Lie group germ, then \(\mathrm{L}(\psi \circ \phi) = \mathrm{L}(\psi) \circ \mathrm{L}(\phi)\) . For \(\phi\) to be an immersion, it is necessary and sufficient that \(\mathrm{L}(\phi)\) be an isomorphism of \(\mathrm{L}(\mathrm{G})\) onto a Lie subalgebra of \(\mathrm{L}(\mathrm{H})\) which admits a topological supplement. In particular, if G is a Lie subgroup germ of H and \(\phi\) is the canonical injection, \(\mathrm{L}(\mathrm{G})\) is identified with a Lie subalgebra of \(\mathrm{L}(\mathrm{H})\) by means of \(\mathrm{L}(\phi)\) . If \((\mathrm{G}_{i})_{i\in I}\) is a finite family of Lie group germs and G is their product, \(\mathrm{L}(\mathrm{G})\) is canonically identified with \(\prod_{i\in I}\mathrm{L}(\mathrm{G}_{i})\) .

18.9. Let G be a Lie group germ, of finite dimension if K is of characteristic \textgreater0. Let F be a complete normable space. Let \(\alpha\) be a differential form of degree k on G with values in F. \(\alpha\) is called left invariant on G if \(\alpha_{g}\) is derived from \(\alpha_{e}\) by the mapping \(h \mapsto gh\) of a neighbourhood of e onto a neighbourhood of g. If \(\alpha\) is left invariant, \(\alpha\) is analytic. The mapping \(\alpha \mapsto \alpha_{e}\) is a bijection of the set of left invariant differential forms of degree k on G with values in F onto the set of continuous alternating k-linear mappings of \(T_{e}(G)\) into F. If \(\alpha_{e} = \operatorname{Id}_{\mathrm{T}_e(\mathrm{G})}\) , \(\alpha\) is called the left canonical differential form of G. The definitions of right invariant differential forms and the right canonical differential form of G are analogous. If \(\omega\) is the left canonical differential form of G, then \(d\omega + [\omega]^{2} = 0\) . Let M be a manifold of class \(C^{r}\) and f a mapping of class \(C^{r}\) of M into G. The left differential of f, denoted by \(f^{-1}.df\) , is the differential form of degree 1 on M with values in L(G) which associates with each vector \(u \in T_{m}(M)\) the element \(f(m)^{-1}\) . \((\mathrm{T}_{m}f)(u)\) . Then \(f^{-1}.df = f^{*}(\omega)\) and \(d\alpha + [\alpha]^{2} = 0\) . If two mappings f and g of M into G have the same left differential and K is of characteristic 0, then \(fg^{-1}\) is locally constant.

